\chapter{Conclusion: the record and the silence}

There are two ways to read these ten books, and the historian who uses them has to decide which one is being done at each sentence.

The first reads them as a source: for prices, dates, buyers, sizes, materials, exhibitions, the sequence of a career. Read that way they are very good and their limits are exactly known. They were kept at the time, by one person, from the cheque; their arithmetic closes; their dates agree with the museums' accession registers where those can be checked, and their measurements agree with the museums' rulers to the stretcher. Where they are wrong they are wrong in the ways a household record is wrong: a dealer's death dated two ways, a hearing dated the 18th and the 19th, a print charged to the wrong plate and corrected on the facing leaf, a name spelt as it sounded. Nothing in them was invented for effect, and the few places where she reconstructed the past from memory, the Armory Show sale at \$200, \$225 and \$250, are the places where the record says three things and can be seen to.

The second reads them as a document: as the work of Josephine Nivison Hopper, who was an actress and then a painter and then, for forty-three years, the keeper of another painter's books, and who wrote into those books, under his drawings, the sentences he would not say. Read that way the ledgers are continuous with the notebooks, and the anecdotes under the paintings, \q{Shirley} and \q{Olga} and \q{Nora} and \q{man night hawk (beek)}, are the trace of a writer who had been given one column to write in and used it. Her descriptions have been quoted by every writer on Hopper since Levin, and usually as evidence of what the pictures are about. They are evidence of what she thought worth recording on the day the picture left the studio, and the notebooks show, with some bitterness, that what she recorded of his was what nobody was recording of hers.

One of the quietest facts in the archive says the same thing from the other side: the only frames priced anywhere in the ten books are the ones Josephine Hopper bought for her own canvases, from a maker in Indiana, and wrote down in the notebook she headed \q{Her book}; his were the gallery's, and they appear in the ledgers as deductions on the dealer's statement.

The silence that the books surround is his. In fifty-four years of Book IV the hand that kept it wrote no sentence that was not a charge, a receipt or a sum. In the three work books his contribution is a drawing, a title and a size. The two sentences of his purpose that the archive preserves are both reported by her: a favourite line of Goethe's, \q{strictly off the record}, and \q{the arrestation of a moment in time - with the most acute realization}, extracted, as she put it, by cruelly pinning him down. Everything else that has been written about what Edward Hopper meant has been written against that silence, and much of it has been written from her sentences without saying so.

\enlargethispage{2\baselineskip}What the edition adds is not a fact but a condition: that every sentence quoted here can be checked against the photograph of the leaf it came from, on one screen, by anybody. The reading is first-pass, unchecked, and dated; its errors will be found; the address of the sheet will not move when they are. That is the right state for a source like this to be in. The books were kept so that someone would read them. The inside cover of Book I says who wrote them and who drew in them, and the notebooks say what it cost. The rest is on the leaves.
